\section{Related Work}
\label{sec:related_work}

\paragraph{Group-equivariant networks.}
\citet{cohen2016group} introduced G-CNNs as a generalisation of standard CNNs in which
the translation group \(\Z^2\) is replaced by larger symmetry groups such as \(p4\) (translations
plus 90\(^\circ\) rotations) and \(p4m\) (translations, 90\(^\circ\) rotations, and reflections). The key
insight is that equivariance to a transformation group, rather than mere invariance, allows
the network to propagate structured representations across layers while still exploiting
weight sharing over the full group orbit. Subsequent work has greatly expanded the scope
of equivariant architectures: harmonic networks \citep{worrall2017harmonic} achieve full
360\(^\circ\) equivariance via circular-harmonic filters; steerable CNNs
\citep{weiler2019general} provide a unified framework for constructing equivariant layers
for arbitrary subgroups of \(\mathrm{E}(2)\); and E(n)-equivariant graph neural networks
\citep{satorras2021egnn} extend equivariance to three-dimensional point clouds. The present
work focuses on the discrete groups \(p4\) and \(p4m\) of \citet{cohen2016group} as the primary
architectural intervention, chosen for their direct relevance to the 90\(^\circ\) rotational
symmetry of histopathology patches.

\paragraph{Equivariant networks for digital pathology.}
\citet{veeling2018rotation} demonstrated that replacing all convolutional layers of a
DenseNet-like architecture with \(p4m\)-equivariant group convolutions improves
tumour-detection performance on the PatchCamelyon benchmark derived from Camelyon16
whole-slide images. Their analysis showed that standard CNNs trained without augmentation
produce unstable, orientation-dependent predictions, while the \gcnn{} produces smooth,
rotation-consistent outputs. Our work extends this line of investigation to the
multi-class tissue-type classification setting and explicitly disentangles the contribution
of architectural equivariance from that of augmentation-based approximation.

\paragraph{Colorectal tissue classification.}
\citet{kather2018100} introduced the NCT-CRC-HE-100K dataset together with a supervised
classification baseline, establishing a nine-class benchmark for computational pathology.
Subsequent work has explored self-supervised pre-training \citep{ciga2022self}, transformer-based
architectures, and multi-magnification approaches on this and related datasets. To our
knowledge, a systematic data-scaling study comparing equivariant and augmented CNNs on
this specific benchmark has not been reported.

\paragraph{Augmentation vs.\ equivariance.}
The relationship between data augmentation and architectural equivariance has been
studied theoretically by \citet{chen2020group}, who showed that training with group
augmentation converges in expectation to an equivariant solution but may require
substantially more data than an architecturally equivariant model. \citet{benton2020learning}
proposed learning the augmentation distribution jointly with the model. Our experimental
design directly tests this relationship empirically under controlled data constraints.